% ---------------------------------------------------------------- Method (P1)
\section{Method: optimal stopping policy gradients}

\paragraph{Stochastic stopping policy.}
A network with parameters $\theta$ reads the history and outputs, at each date,
$\varphi^\theta_j=\mathbb P(\text{stop at } j\mid\text{not stopped before},\ \mathcal F_j)$,
with $\varphi_H=1$. Because stopping does not change the path, the distribution of the
stopping time on a given path is exact~\cite{ospg2023}:
\begin{equation}
\psi^\theta_j=\mathbb P(\tau=j)=\varphi^\theta_j\prod_{n<j}\big(1-\varphi^\theta_n\big),\qquad
\psi^\theta_H=\prod_{n<H}\big(1-\varphi^\theta_n\big).
\end{equation}
The expected reward $J(\theta)=\mathbb E\big[\sum_{j}\psi^\theta_j g_j\big]$ can therefore be
evaluated on pre-simulated paths, without rollouts and without backward recursion.

\paragraph{Gradient and loss.}
By the log-derivative trick,
\begin{equation}
\nabla_\theta J=\mathbb E\Big[\sum_{j=0}^{H}(g_j-b)\,\psi^\theta_j\,\nabla_\theta\log\psi^\theta_j\Big],
\end{equation}
for any constant $b$, since $\sum_j\psi_j=1$ makes the $b$ term vanish. We minimise the
temporal loss $\mathcal L=-\frac1N\sum_{i,j}\mathrm{sg}\big[\psi_{ij}(g_{ij}-b)\big]\log\psi_{ij}$,
where $\mathrm{sg}$ stops gradients, whose gradient is exactly $-\nabla_\theta J$ (checked
numerically in our tests). The update raises $\varphi_j$ when the ratio of the immediate
value $\psi_jg_j$ to the continuation value $\sum_{n>j}\psi_ng_n$ exceeds the current odds
$\varphi_j/(1-\varphi_j)$, an implicit Snell-envelope comparison at every date.

\paragraph{Architecture and training.}
The recurrent policy (OSPG-GRU) is a single-layer GRU with 20 hidden units followed by a
linear layer and a sigmoid; the memoryless policy (OSPG-MLP) applies a two-layer
network with 20 units to each date's inputs, with batch normalisation of its input and of
each hidden layer's output, as in the paper; at decision time the normalisation uses the
statistics from training, so a path's decision never depends on the other paths. Inputs at date $j$ are the standardised
features, the current reward and $j/H$, all known at date~$j$. We train with Adam
(learning rate $10^{-3}$), keep the parameters with the best validation value, and stop
at the first date with $\varphi_j\ge\frac12$ at test time. Three implementation choices
were needed. The output bias starts at $-2$, because a policy started near
$\varphi=\frac12$ converged to ``stop immediately'' on the in-the-money put. A small
penalty $10^{-3}\,\overline{z^2}$ on the pre-sigmoid outputs keeps $\varphi$ away from $0$:
the policy gradient is proportional to $\varphi$, and without the penalty the GRU policy
stalled at ``never stop'' on the 50-asset max-call. Finally, date~$0$, where all simulated
paths coincide, is decided by comparing $g_0$ with the value achieved by the trained
policy, as for the baselines.

\paragraph{Baselines.}
Longstaff--Schwartz regression on $[1,x,x^2]$ of the standardised features and current
reward (in-the-money paths only for options); Deep Optimal Stopping with one network per
date, warm-started from the next date; and fitted Q-iteration with one network (MLP or
GRU) for all dates~\cite{ospg2023}. The baseline networks are run without batch
normalisation: on the American put it lowered FQI-MLP from $4.480$ and $4.462$ to $4.426$
and $4.442$ (two seeds), so we kept the stronger configuration for the baselines, which
can only make the comparison harder for OSPG. Every method sees the same features, and
all values are measured out of sample, so they are lower bounds on the true price; the
binomial tree and the finite-difference solver of Section~2 give the exact prices they
are measured against.
